\clearpage
\section{FDT}
\begin{colbox}
	The fluctuation-dissipation theorem (FDT) related the correlation function $C(t)$ with the response function $\chi(t)$ according to
	\begin{equation}\label{diffC}
		\diffp{C}{t} = -k_BT\chi(t) \quad t>0
	\end{equation}
	we define the fourier transformations
	\begin{align}
		\tilde C(\omega) &= \int\limits_{-\infty}^\infty dt C(t)e^{i\omega t}\\
		\tilde \chi(\omega) &= \int\limits_0^\infty dt \chi(t)e^{i\omega t} \equiv \tilde \chi'(\omega)+i\tilde\chi''(\omega)
	\end{align}
	where $\tilde\chi'(\omega)$ and $\tilde\chi''(\omega)$ are the real and imaginary parts of the transformed response function, respectively. Show that $C(t)$ is an even function and derive FDT in the frequency domain,
	\begin{equation}
		\tilde C(\omega) = \frac{2k_BT}{\omega}\tilde\chi''(\omega)
	\end{equation}
\end{colbox}
In the lecture we saw that $C$ does not depend on the explicit times but rather only on the time that has passed, therefore we can rewrite
\begin{equation}
	C(t) = \langle\delta x(0)\delta x(t)\rangle \rightarrow C(t) = \langle\delta x(-t/2)\delta x(t/2)\rangle
\end{equation}
which we can immediately notice that, assuming the system is invariant to time reversal,
\begin{equation}
	C(-t) = \langle\delta x(t/2)\delta x(-t/2)\rangle = C(t)
\end{equation}
therefore $C$ is even. Next, we look at
\begin{equation}
	\diffp{C}{t} = -k_BT\chi(t) \quad t>0
\end{equation}
and notice that since $C$ is even w.r.t time, its first time derivative must be odd. This means $\chi$ must be odd as well and only depend on the absolute amount of time between t and 0:
\begin{equation}
	\diffp{C}{t} = -k_BT\chi(\lvert t \rvert)\text{sgn}(t)
\end{equation}
we'll take the fourier transform of the derivative
\begin{align}
	\int\limits_{-\infty}^\infty dt \diffp{C}{t}e^{i\omega t} &= -k_BT\int\limits_{-\infty}^\infty dt \chi(\lvert t \rvert)\text{sgn}(t)e^{i\omega t}\\
	&= -k_BT\ins{\int\limits_0^\infty dt \chi(\lvert t \rvert)e^{i\omega t}-\int\limits_{-\infty}^0 dt \chi(\lvert t \rvert)e^{i\omega t}}\\
	&= -k_BT\ins{\tilde\chi(\omega)-\tilde\chi(-\omega)}
\end{align}
since $\chi$ is real, we can use the identity $\tilde\chi(-\omega) = \tilde\chi^*(\omega)$ and find
\begin{align}
	\tilde\chi(\omega)-\tilde\chi(-\omega) &= \tilde\chi(\omega)-\tilde\chi^*(\omega)\\
	&= 2i\Im(\tilde\chi(\omega))\\
	&= 2i\tilde\chi''(\omega)
\end{align}
as for the other side of the equation we can use another fourier transform identity to find
\begin{equation}
	\int\limits_{-\infty}^\infty dt \diffp{C}{t}e^{i\omega t} = i\omega\tilde C(\omega)
\end{equation}
therefore
\begin{equation}
	i\omega\tilde C(\omega) = 2ik_BT\tilde\chi''(\omega)
\end{equation}
cancelling out terms we find
\begin{equation}
	\tilde C(\omega) = \frac{2k_BT}{\omega}\tilde\chi''(\omega)
\end{equation}
as required.